\section{Conventions}\label{sec:conventions}

Many binary choices enter the theory of modular curves: rows or columns,
left or right actions, the direction of a level structure, the lattice
attached to a point $\tau$ of the upper half plane. They influence each
other, and the literature does not agree on them. In this section we make
all choices explicit, starting from objects with an intrinsic meaning,
and we compare with the literature at the end.

\subsection{Matrices and level structures}

\begin{convention}[Matrices act on columns]\label{conv:columns}
For a commutative ring $R$ the elements of $\StdModule{R}{2}$ are column
vectors, $\StdBasis{1} = \transpose{(1,0)}$ and
$\StdBasis{2} = \transpose{(0,1)}$, and $A \in \MatRing{2}{R}$ acts on
$\StdModule{R}{2}$ by $v \mapsto Av$. In particular the linear maps
attached to $A$ and $B$ compose as $v \mapsto ABv$, and every group of
matrices acts on $\StdModule{R}{2}$ on the left.
\end{convention}

\begin{convention}[Level structures]\label{conv:level}
Let $N \geq 1$ and let $E$ be an elliptic curve over a field $k$ of
characteristic not dividing $N$. A \emph{full level $N$ structure} on $E$
is a group isomorphism
$\LevelStructure \colon \StdModule{(\ResidueRing{N})}{2} \to \Torsion{E}{N}(\overline{k})$,
as in \cite{KatzMazur1985}. The group $\GL{2}{\ResidueRing{N}}$ acts on the
set of full level $N$ structures on the right by precomposition,
$\LevelStructure \cdot h = \LevelStructure \circ h$. For a subgroup
$H \subseteq \GL{2}{\ResidueRing{N}}$ an \emph{$H$-level structure} is an
orbit $\LevelClass{\LevelStructure}{H}$.
\end{convention}

\begin{remark}
Deligne and Rapoport use the opposite direction
$\Torsion{E}{N} \to \StdModule{(\ResidueRing{N})}{2}$, on which
$\GL{2}{\ResidueRing{N}}$ acts on the left
\cite[IV.2.3, IV.3.1]{DeligneRapoport1973}. Passing from one to the other
is $\LevelStructure \mapsto \LevelStructure^{-1}$, which also moves the
action of $H$ to the other side.
\end{remark}

\begin{example}[The Borel subgroup]\label{ex:borel-moduli}
Let $\BorelMod{N} \subseteq \GL{2}{\ResidueRing{N}}$ be the group of upper
triangular matrices. For $h = \smallMat{a}{b}{0}{d} \in \BorelMod{N}$ one
has $(\LevelStructure \circ h)(\StdBasis{1}) = a\,\LevelStructure(\StdBasis{1})$.
Hence a $\BorelMod{N}$-level structure on $E$ is the same thing as the
cyclic subgroup $\langle \LevelStructure(\StdBasis{1}) \rangle$ of order
$N$. Similarly, for the group of matrices
$\smallMat{1}{b}{0}{d}$ an $H$-level structure is the same thing as the
point $\LevelStructure(\StdBasis{1})$ of order $N$.
\end{example}

\subsection{The lattice attached to a point of the slit plane}

\begin{definition}[Period vector, lattice, level structure of $\tau$]
\label{def:lattice-of-tau}
Let $\tau \in \SlitPlane$. The \emph{period vector} of $\tau$ is
$\PeriodVector{\tau} := \transpose{(\tau, 1)} \in \StdModule{\CC}{2}$.
Define the $\RR$-linear map
\[
  \LatticeBasis{\tau} \colon \StdModule{\RR}{2} \to \CC, \qquad
  v \mapsto \DetPair{\PeriodVector{\tau}}{v},
\]
the determinant of the $2\times 2$ complex matrix with columns
$\PeriodVector{\tau}$ and $v$. Put
$\LatticeOf{\tau} := \LatticeBasis{\tau}(\StdModule{\ZZ}{2})$ and
$\EllCurve{\tau} := \CC/\LatticeOf{\tau}$. For $N \geq 1$ let
\[
  \StdLevel{\tau}{N} \colon \StdModule{(\ResidueRing{N})}{2} \to
  \Torsion{\EllCurve{\tau}}{N}, \qquad
  v \mapsto \tfrac{1}{N}\LatticeBasis{\tau}(v) \bmod \LatticeOf{\tau}.
\]
\end{definition}

The point of using the determinant pairing, rather than the dot product
$v \mapsto \transpose{\PeriodVector{\tau}} v$, is that it is equivariant
for $\GL{2}{\RR}$ (\cref{lem:det-pairing-equivariance}); the dot product is
only equivariant for the orthogonal group, which is where the transposes
in other conventions come from (\cref{rem:literature}). Intrinsically,
$\StdModule{\ZZ}{2}$ plays the role of $\Homology{\EllCurve{\tau}}{\ZZ}$, the
determinant is its Riemann form, $\PeriodVector{\tau}$ spans its Hodge
filtration, and $\LatticeBasis{\tau}$ is the period map of a normalised
differential; see \cref{app:abelian}, in particular
\cref{prop:frame-tau}.

\begin{lemma}[Coordinates of the lattice of $\tau$]\label{lem:lattice-coordinates}
\slogan{The lattice of $\tau$ has basis $-1, \tau$.}
Let $\tau \in \SlitPlane$ and $N \geq 1$. Then
$\LatticeBasis{\tau}(\StdBasis{1}) = -1$,
$\LatticeBasis{\tau}(\StdBasis{2}) = \tau$,
$\LatticeOf{\tau} = \ZZ + \ZZ\tau$ is a lattice in $\CC$, and
$\StdLevel{\tau}{N}$ is a well defined group isomorphism. Moreover
$\LatticeBasis{\tau}(v) = \transpose{\PeriodVector{\tau}}\,\Sform\, v$ with
$\Sform = \smallMat{0}{1}{-1}{0}$.
\end{lemma}

\begin{proof}
For $v = \transpose{(x,y)}$ one has
$\DetPair{\PeriodVector{\tau}}{v} = \tau y - x =
\transpose{\PeriodVector{\tau}} \Sform v$. Since $\tau \notin \RR$, the
numbers $-1, \tau$ are $\RR$-linearly independent, so $\LatticeOf{\tau}$ is
a lattice, $\LatticeBasis{\tau}$ restricts to an isomorphism
$\StdModule{\ZZ}{2} \to \LatticeOf{\tau}$, and hence induces an isomorphism
$\StdModule{(\ResidueRing{N})}{2} \cong \tfrac{1}{N}\LatticeOf{\tau}/\LatticeOf{\tau}
= \Torsion{\EllCurve{\tau}}{N}$.
\end{proof}

\begin{lemma}[Inverse transpose is conjugation by $\Sform$]\label{lem:S-identity}
\slogan{Transposes in $\GL{2}{}$ are inner, up to the determinant.}
Let $R$ be a commutative ring and $\gamma \in \GL{2}{R}$. Then
$\Sform \gamma \Sform^{-1} = \det(\gamma)\,(\transpose{\gamma})^{-1}$,
where $\Sform = \smallMat{0}{1}{-1}{0}$.
\end{lemma}

\begin{proof}
For $\gamma = \smallMat{a}{b}{c}{d}$ both sides equal
$\smallMat{d}{-c}{-b}{a}$.
\end{proof}

\begin{lemma}[M\"obius transformations are the left action on lines]
\label{lem:mobius}
\slogan{$\gamma$ moves the period vector of $\tau$ to a multiple of the
period vector of $\gamma\tau$.}
Let $\gamma = \smallMat{a}{b}{c}{d} \in \GL{2}{\RR}$ and
$\tau \in \SlitPlane$. Then $c\tau + d \neq 0$,
$\MobiusAction{\gamma}{\tau} := (a\tau+b)/(c\tau+d) \in \SlitPlane$, and
$\gamma \PeriodVector{\tau} = (c\tau + d)\,\PeriodVector{\gamma\tau}$.
This defines a left action of $\GL{2}{\RR}$ on $\SlitPlane$, which
preserves $\UpperHalfPlane$ for elements of positive determinant.
\end{lemma}

\begin{proof}
Standard: $\gamma\PeriodVector{\tau} = \transpose{(a\tau+b, c\tau+d)}$, and
$\mathrm{Im}(\gamma\tau) = \det(\gamma)\,\mathrm{Im}(\tau)/|c\tau+d|^{2}$.
\end{proof}

\begin{lemma}[Equivariance of the determinant pairing]
\label{lem:det-pairing-equivariance}
\slogan{The lattice of $\gamma\tau$ is the lattice of $\tau$ transported
by $\gamma$.}
Let $\gamma = \smallMat{a}{b}{c}{d} \in \GL{2}{\RR}$, $\tau \in \SlitPlane$
and $v \in \StdModule{\RR}{2}$. Then
\[
  \LatticeBasis{\gamma\tau}(\gamma v) =
  \det(\gamma)\,(c\tau + d)^{-1}\,\LatticeBasis{\tau}(v).
\]
In particular, for $\gamma \in \GL{2}{\ZZ}$ the map
$z \mapsto \det(\gamma)(c\tau+d)^{-1} z$ is an isomorphism
$\EllCurve{\tau} \to \EllCurve{\gamma\tau}$ which transforms
$\StdLevel{\tau}{N}$ into $\StdLevel{\gamma\tau}{N} \circ \gamma$ for every
$N \geq 1$.
\end{lemma}

\begin{proof}
By the left action on lines (\cref{lem:mobius}),
$\LatticeBasis{\gamma\tau}(\gamma v)
= (c\tau+d)^{-1}\DetPair{\gamma\PeriodVector{\tau}}{\gamma v}
= (c\tau+d)^{-1}\det(\gamma)\DetPair{\PeriodVector{\tau}}{v}$.
For $\gamma \in \GL{2}{\ZZ}$ one has $\gamma\StdModule{\ZZ}{2} = \StdModule{\ZZ}{2}$,
so the displayed identity says that multiplication by
$\det(\gamma)(c\tau+d)^{-1}$ maps $\LatticeOf{\tau}$ onto
$\LatticeOf{\gamma\tau}$, compatibly with the maps $\LatticeBasis{\tau}$ and
$\LatticeBasis{\gamma\tau} \circ \gamma$.
\end{proof}

\begin{lemma}[Weil determinant of the standard level structure]
\label{lem:weil-determinant}
Normalise the Weil pairing on $\EllCurve{\tau}$, $\tau \in \UpperHalfPlane$,
such that $\WeilPairing{N}(\tau/N, 1/N) = \exp(2\pi i/N)$, as in
\cite{DiamondShurman2005}. Then
$\WeilPairing{N}(\StdLevel{\tau}{N}(\StdBasis{1}), \StdLevel{\tau}{N}(\StdBasis{2}))
= \exp(2\pi i/N)$.
\end{lemma}

\begin{proof}
By the coordinates of the lattice of $\tau$ (\cref{lem:lattice-coordinates})
the left hand side is $\WeilPairing{N}(-1/N, \tau/N)$, which equals
$\WeilPairing{N}(\tau/N, 1/N)$ because the pairing is bilinear and
alternating.
\end{proof}

\todo[inline]{Double check the normalisation of the Weil pairing in
\cite{DiamondShurman2005} against \cite{KatzMazur1985} and
\cite{DeligneRapoport1973}.}

\subsection{Complex uniformisation}

\begin{proposition}[Complex points of $\OpenModCurve{H}$]
\label{prop:complex-uniformisation}
\slogan{Elliptic curves with $H$-level structure are pairs $(\tau, g)$ up
to $\GL{2}{\ZZ}$ on the left and $H$ on the right.}
Let $N \geq 1$ and $H \subseteq \GL{2}{\ResidueRing{N}}$ a subgroup. Let
$\GL{2}{\ZZ}$ act on $(\SlitPlane) \times \GL{2}{\ResidueRing{N}}$ on the
left by $\gamma\cdot(\tau, g) = (\gamma\tau, \overline{\gamma} g)$, where
$\overline{\gamma}$ is the reduction of $\gamma$ modulo $N$, and let $H$
act on the right on the second factor. Then
\begin{gather*}
  \Phi \colon \GL{2}{\ZZ}\backslash\bigl((\SlitPlane)\times\GL{2}{\ResidueRing{N}}\bigr)/H
  \longrightarrow
  \{(E, \LevelClass{\LevelStructure}{H})\}/\cong,\\
  (\tau, g) \mapsto (\EllCurve{\tau}, \LevelClass{(\StdLevel{\tau}{N}\circ g)}{H}),
\end{gather*}
where $E$ runs over elliptic curves over $\CC$, is a well defined
bijection. It is compatible with the right actions of
$\GL{2}{\ResidueRing{N}}$ when $H = \{1\}$.
\end{proposition}

\begin{proof}
Well defined: by the equivariance of the determinant pairing
(\cref{lem:det-pairing-equivariance}),
$(\EllCurve{\gamma\tau}, \StdLevel{\gamma\tau}{N}\circ\overline{\gamma} g)
\cong (\EllCurve{\tau}, \StdLevel{\tau}{N}\circ g)$ for $\gamma \in \GL{2}{\ZZ}$.

Surjective: write $E = \CC/L$ and choose a basis
$\omega \colon \StdModule{\ZZ}{2} \to L$. Put $\lambda := -\omega(\StdBasis{1})$
and $\tau := -\omega(\StdBasis{2})/\omega(\StdBasis{1}) \in \SlitPlane$.
Then $\omega = \lambda\LatticeBasis{\tau}$ by the coordinates of the lattice
of $\tau$ (\cref{lem:lattice-coordinates}), so $z \mapsto \lambda^{-1}z$ is
an isomorphism $E \to \EllCurve{\tau}$, and it transforms a level structure
$\LevelStructure$ into $\StdLevel{\tau}{N}\circ g$ with
$g := \StdLevel{\tau}{N}^{-1}\circ\lambda^{-1}\LevelStructure$.

Injective: suppose $z \mapsto \mu z$ is an isomorphism
$\EllCurve{\tau} \to \EllCurve{\tau'}$ transforming
$\StdLevel{\tau}{N} \circ g$ into an element of
$\LevelClass{(\StdLevel{\tau'}{N}\circ g')}{H}$. Both $\mu\LatticeBasis{\tau}$
and $\LatticeBasis{\tau'}$ restrict to bases of $\LatticeOf{\tau'}$, so
$\mu\LatticeBasis{\tau} = \LatticeBasis{\tau'}\circ\gamma$ for a unique
$\gamma \in \GL{2}{\ZZ}$. Comparing with
$\LatticeBasis{\gamma\tau}\circ\gamma = \det(\gamma)(c\tau+d)^{-1}\LatticeBasis{\tau}$
shows that $\LatticeBasis{\tau'}$ and $\LatticeBasis{\gamma\tau}$ are
proportional, hence $\tau' = \gamma\tau$ (the ratio
$-\LatticeBasis{\tau}(\StdBasis{2})/\LatticeBasis{\tau}(\StdBasis{1})$
recovers $\tau$), and that $\mu\StdLevel{\tau}{N} = \StdLevel{\tau'}{N}\circ\overline{\gamma}$.
Therefore $\overline{\gamma} g \in g' H$, i.e.\ $(\tau', g') \in
\gamma\cdot(\tau, g)\cdot H$.
\end{proof}

\begin{corollary}[Classical uniformisation without transposes]
\label{cor:classical-uniformisation}
\slogan{$\OpenModCurve{H}(\CC) = \GammaOfLevel{H}\backslash\UpperHalfPlane$
with $\GammaOfLevel{H} = H \cap \SL{2}{\ZZ}$.}
Let $N \geq 1$ and let $H \subseteq \GL{2}{\ResidueRing{N}}$ be a subgroup
with $\det(H) = \unit{(\ResidueRing{N})}$. Let
$\GammaOfLevel{H} := \{\gamma \in \SL{2}{\ZZ} : \overline{\gamma} \in H\}$.
Then $\tau \mapsto (\tau, 1)$ induces a bijection
\[
  \GammaOfLevel{H}\backslash\UpperHalfPlane \longrightarrow
  \GL{2}{\ZZ}\backslash\bigl((\SlitPlane)\times\GL{2}{\ResidueRing{N}}\bigr)/H .
\]
\end{corollary}

\begin{proof}
Surjective: given $(\tau, g)$, replace it by
$\Diagonal{1}{-1}\cdot(\tau, g)$ if necessary to get
$\tau \in \UpperHalfPlane$. Choose $h \in H$ with $\det(gh) = 1$ and a lift
$\sigma \in \SL{2}{\ZZ}$ of $gh$ (reduction $\SL{2}{\ZZ} \to \SL{2}{\ResidueRing{N}}$
is surjective). Then $(\tau, g)H = \sigma\cdot(\sigma^{-1}\tau, 1)H$.
Injective: if $\gamma\cdot(\tau, 1) = (\tau', h)$ with
$\tau, \tau' \in \UpperHalfPlane$ and $h \in H$, then $\det\gamma = 1$
because $\gamma$ preserves $\UpperHalfPlane$, and
$\overline{\gamma} = h \in H$, so $\gamma \in \GammaOfLevel{H}$.
\end{proof}

\begin{example}[The groups $\GammaZero{N}$, $\GammaOne{N}$, $\GammaFull{N}$]
\label{ex:classical-groups}
For $H = \BorelMod{N}$ one gets
$\GammaOfLevel{H} = \{\smallMat{a}{b}{c}{d} \in \SL{2}{\ZZ} : N \mid c\}$, and
$\tau$ corresponds to $(\EllCurve{\tau}, \langle 1/N\rangle)$ by
the moduli description of $\BorelMod{N}$ (\cref{ex:borel-moduli}), since
$\StdLevel{\tau}{N}(\StdBasis{1}) = -1/N$. For
$H = \{\smallMat{1}{b}{0}{d}\}$ one gets the group of matrices with
$c \equiv 0$ and $a \equiv d \equiv 1 \bmod N$, and $\tau$ corresponds to
$(\EllCurve{\tau}, -1/N) \cong (\EllCurve{\tau}, 1/N)$. For $H = \{1\}$ one
gets the principal congruence subgroup, and $\tau$ corresponds to
$(\EllCurve{\tau}, (-1/N, \tau/N))$. These agree with
\cite[Thm.~1.5.1]{DiamondShurman2005} and \cite[\S 7.2]{DiamondIm1995}.
\end{example}

\subsection{The adelic description}

\begin{lemma}[Class number one for $\GL{2}{}$]\label{lem:gl2-class-number-one}
\slogan{Every lattice in $\FiniteAdeles{\QQ}^{2}$ is rational.}
One has $\GL{2}{\FiniteAdeles{\QQ}} = \GL{2}{\QQ}\,\GL{2}{\ProfiniteIntegers}$
and $\GL{2}{\QQ} \cap \GL{2}{\ProfiniteIntegers} = \GL{2}{\ZZ}$ inside
$\GL{2}{\FiniteAdeles{\QQ}}$.
\end{lemma}

\begin{proof}
For $g \in \GL{2}{\FiniteAdeles{\QQ}}$ the set
$L := g\StdModule{\ProfiniteIntegers}{2} \cap \StdModule{\QQ}{2}$ is a
lattice in $\StdModule{\QQ}{2}$ whose closure is
$g\StdModule{\ProfiniteIntegers}{2}$; writing $L = \gamma\StdModule{\ZZ}{2}$
with $\gamma \in \GL{2}{\QQ}$ gives $\gamma^{-1}g \in \GL{2}{\ProfiniteIntegers}$.
The second statement holds because a rational matrix that is integral at
all primes, with determinant a unit at all primes, lies in $\GL{2}{\ZZ}$.
\end{proof}

\begin{definition}[Adelic double coset]\label{def:adelic-double-coset}
For an open subgroup $\mathcal{K} \subseteq \GL{2}{\FiniteAdeles{\QQ}}$ put
\[
  \OpenModCurve{\mathcal{K}}(\CC) := \GL{2}{\QQ}\backslash\bigl((\SlitPlane)\times
  \GL{2}{\FiniteAdeles{\QQ}}\bigr)/\mathcal{K},
\]
where $\GL{2}{\QQ}$ acts by $\gamma\cdot(\tau, g) = (\gamma\tau, \gamma g)$
(diagonal left multiplication) and $\mathcal{K}$ acts by right
multiplication on the second factor.
\end{definition}

\begin{proposition}[Finite level versus adelic level]\label{prop:finite-vs-adelic}
\slogan{The adelic double coset with plain left multiplication recovers
the finite level description.}
Let $N \geq 1$, $H \subseteq \GL{2}{\ResidueRing{N}}$ a subgroup and
$\mathcal{K} \subseteq \GL{2}{\ProfiniteIntegers}$ its preimage. Then the
inclusion $\GL{2}{\ProfiniteIntegers} \subseteq \GL{2}{\FiniteAdeles{\QQ}}$
and reduction modulo $N$ induce bijections
\[
  \GL{2}{\ZZ}\backslash\bigl((\SlitPlane)\times\GL{2}{\ResidueRing{N}}\bigr)/H
  \cong
  \GL{2}{\ZZ}\backslash\bigl((\SlitPlane)\times\GL{2}{\ProfiniteIntegers}\bigr)/\mathcal{K}
  \cong \OpenModCurve{\mathcal{K}}(\CC).
\]
\end{proposition}

\begin{proof}
The first bijection holds because
$\GL{2}{\ProfiniteIntegers}/\mathcal{K} = \GL{2}{\ResidueRing{N}}/H$. For the
second, surjectivity follows from class number one for $\GL{2}{}$
(\cref{lem:gl2-class-number-one}), and if
$\gamma k u = k'$ with $\gamma \in \GL{2}{\QQ}$, $k, k' \in \GL{2}{\ProfiniteIntegers}$
and $u \in \mathcal{K}$ then
$\gamma \in \GL{2}{\QQ}\cap\GL{2}{\ProfiniteIntegers} = \GL{2}{\ZZ}$.
\end{proof}

\begin{remark}[Comparison with the literature]\label{rem:literature}
\begin{enumerate}
\item Diamond--Im \cite[\S 11.1, Rem.~11.1.1]{DiamondIm1995}, Shimura
  \cite{Shimura1971}, Deligne \cite{DeligneTravaux1971,DeligneVarietes1979}
  and Milne \cite{Milne2005} all write
  $G(\QQ)\backslash X \times G(\FiniteAdeles{\QQ})/K$ with plain left
  multiplication by $G(\QQ)$, as in \cref{def:adelic-double-coset}.
  Diamond--Im define $U_0(N)$ by $c \equiv 0 \bmod N$, in agreement with
  \cref{ex:classical-groups}. Their $U_1(N)$ is defined by $c \equiv 0$ and
  $d \equiv 1$ instead of $a \equiv 1$; on $\SL{2}{}$ the two conditions
  agree.
  \todo[inline]{Check whether the $d \equiv 1$ version corresponds to the
  $\mu_N \hookrightarrow E$ moduli problem $\mathcal{Y}_\mu(N)$ of
  \cite[Variant 8.2.2]{DiamondIm1995} in our conventions.}
\item Diamond--Shurman \cite{DiamondShurman2005} use the point $1/N$ and
  the basis $(\tau/N, 1/N)$ of $\Torsion{\EllCurve{\tau}}{N}$; this agrees
  with \cref{ex:classical-groups} after the change of basis $\Sform$.
\item Rouse--Sutherland--Zureick-Brown \cite[\S 2]{RSZB2022} use column
  vectors with a left action, upper triangular $\BorelMod{N}$, level
  structures $\Torsion{E}{N} \to \StdModule{(\ResidueRing{N})}{2}$ with $H$
  acting by postcomposition, and $\GammaOfLevel{H} = \pm H \cap \SL{2}{}$
  without transpose. Up to $\LevelStructure \mapsto \LevelStructure^{-1}$
  this agrees with \cref{conv:level} and
  \cref{cor:classical-uniformisation}. They note that
  \cite{RouseZureickBrown2015} uses right actions.
\item Zywina \cite[Remark 3.6]{ZywinaExplicit2022} defines
  $\GammaOfLevel{G} = G \cap \SL{2}{\ZZ}$ without transpose, but works with
  the dual representation $\sigma \mapsto \transpose{\rho_E(\sigma^{-1})}$,
  and warns that $\ModCurve{G}$ in other sources may be his
  $\ModCurve{\transpose{G}}$.
\item Deligne--Rapoport \cite[IV.5.1--5.3]{DeligneRapoport1973} use a basis
  $\beta \colon \StdModule{\ZZ}{2} \to H_1(E,\ZZ)$ with
  $z(\beta) = \beta(\StdBasis{1})/\beta(\StdBasis{2})$, a right action of
  $\GL{2}{\ZZ}$ on bases written with transposed matrix entries, and the
  double coset
  $K\backslash X^{\pm}\times\GL{2}{\FiniteAdeles{\QQ}}/\GL{2}{\QQ}$. For our
  basis $\LatticeBasis{\tau}$ one has $z = -1/\tau = \MobiusAction{\Sform}{\tau}$, and
  their double coset is identified with ours by $g \mapsto g^{-1}$,
  $\LevelStructure \mapsto \LevelStructure^{-1}$.
\item Conway \cite{Conway1996} uses row vectors with a right action; his
  statements translate to ours by transposition.
\item In all cases the transposes that appear are explained by
  inverse transpose being conjugation by $\Sform$
  (\cref{lem:S-identity}), i.e.\ by the difference between the dot product
  and the determinant pairing.
\end{enumerate}
\end{remark}
